% META: {"id": "1SPE-TRIGO-CO-037", "chapitre": "1SPE-TRIGONOMETRIE", "type_objet": "corrige", "exercice_ref": "1SPE-TRIGO-EX-037", "capacites_codes": ["C4"], "sources_inspiration": [], "status": "generated"}
% BEGIN-VERIFY
% from sympy import *
% assert cos(3*pi/4) == -sqrt(2)/2
% assert cos(5*pi/4) == -sqrt(2)/2
% assert sin(pi/3) == sqrt(3)/2
% assert sin(2*pi/3) == sqrt(3)/2
% END-VERIFY
\begin{corrige}{1SPE-TRIGO-EX-037}

\textbf{1.} $\cos x \leqslant -\dfrac{\sqrt{2}}{2}$ : $\cos x = -\dfrac{\sqrt{2}}{2}$ en $x = \dfrac{3\pi}{4}$ et $x = \dfrac{5\pi}{4}$. Le cosinus est inferieur a $-\dfrac{\sqrt{2}}{2}$ pour $x \in \left[\dfrac{3\pi}{4};\dfrac{5\pi}{4}\right]$.

\textbf{2.} $\sin x \geqslant \dfrac{\sqrt{3}}{2}$ : $\sin x = \dfrac{\sqrt{3}}{2}$ en $x = \dfrac{\pi}{3}$ et $x = \dfrac{2\pi}{3}$. Le sinus est superieur ou egal a $\dfrac{\sqrt{3}}{2}$ pour $x \in \left[\dfrac{\pi}{3};\dfrac{2\pi}{3}\right]$.

\end{corrige}
